% 核心观点：示范动作误差不是任务表现本身；任务改善取决于部署访问情形中的动作后果。
% 叙述逻辑：B1定义动作后果与性能差分；B2给出覆盖、敏感性条件下的动作误差界；
% B3用表示风险迁移的核心公式提出关系信息与执行更新两个问题；详细证明留附录，动作MSE诊断留D及附录。
% C、D继续研究表示与反馈的条件性作用，不预先声称语义降低C或反馈降低L。
% 边界：覆盖界只刻画动作风险迁移所需的条件；完整动作历史可能导致结构性支持失配。
% 不假定专家在离线数据之外可查询，也不由固定示范推出覆盖或随N变化的统计保证。
\subsection{From Imitation Error to Task Performance}
\label{sec:actionanalysis}
% 动作误差能否约束任务失败，取决于误差发生的情形及其后续任务后果；以下建立这一联系，并分析表示空间风险迁移的机会与限制。
Whether action error bounds task failure depends on where deviations occur and their subsequent task consequences.
We establish this connection and examine the opportunities and limitations of risk transfer through representations.

% \subsubsection{动作对任务的影响}
\subsubsection{Effects of Actions on Task Outcomes}
\label{sec:actionconsequences}
% 在同一执行情形下比较不同动作，随后均沿同一参考策略 $\pi_0$ 继续执行，
To compare actions in the same execution situation, we follow each action with the same reference policy $\pi_0$.
% 即可用最终失败概率的差异衡量当前动作对任务的影响。
The resulting difference in failure probability measures the effect of the current action on the task.
% 为保留任务时序与策略历史，引入分析用扩展状态 $x_t=(s_{0:t},\mathcal H_t)$。
To retain task timing and policy history, we introduce the augmented analysis state $x_t=(s_{0:t},\mathcal H_t)$.
% 后续分析限于固定的确定性策略，并简写 $\pi(x_t)=\pi(\mathcal H_t)$；
The following analysis considers fixed deterministic policies and uses the shorthand $\pi(x_t)=\pi(\mathcal H_t)$;
% 策略实际仅使用可获得的历史 $\mathcal H_t$。
the policy itself uses only the available history $\mathcal H_t$.

% 在控制步 $t$ 的执行情形 $x$ 下先执行动作 $a$、随后遵循 $\pi_0$ 时，
Let $Q_t^0(x,a)$ denote the conditional probability of task failure when action $a$ is executed
% 任务失败的条件概率记为 $Q_t^0(x,a)$。
in situation $x$ at control step $t$ and $\pi_0$ is followed thereafter.
% 执行参考动作时的失败概率为 $V_t^0(x)=Q_t^0(x,\pi_0(x))$，
The failure probability under the reference action is $V_t^0(x)=Q_t^0(x,\pi_0(x))$.
% 当前动作相对参考动作造成的失败概率差为
The change in failure probability relative to that action is
$A_t^0(x,a)=Q_t^0(x,a)-V_t^0(x)$.
% $A_t^0(x,a)>0$ 表示该动作增加失败概率，$A_t^0(x,a)<0$ 表示降低失败概率。
Thus, $A_t^0(x,a)>0$ indicates an increase in failure probability, whereas $A_t^0(x,a)<0$ indicates a decrease.

% \begin{assumption}[共同的比较过程]
\begin{assumption}[Common comparison process]
\label{ass:commonprocess}
% 比较策略共享初始分布、受控转移和任务完成函数。
The compared policies share the initial distribution, controlled transitions, and task-completion function.
% 包含运行记录时，共同转移也包括计算与记录过程，而不只是相同的机器人动力学。
When runtime records are included, the common transitions also cover the computation and recording process, not only the robot dynamics.
% 参考策略 $\pi_0$ 在比较所涉及的全部历史上有定义；
The reference policy $\pi_0$ is defined on every history involved in the comparison.
% 动作规则和上述条件期望可测，执行动作属于 $\mathcal A$。
The action rules and the conditional expectations above are measurable, and executed actions belong to $\mathcal A$.
\end{assumption}

% 以 $F(\pi)=1-J(\pi)$ 表示策略 $\pi$ 的任务失败概率。
Let $F(\pi)=1-J(\pi)$ denote the task-failure probability of policy $\pi$.
% 待评价策略诱导的时刻--历史分布 $\mu_\pi^t$ 描述其在控制步 $t$ 实际访问的 $x_t$，
The time--history distribution $\mu_\pi^t$ describes the augmented histories $x_t$ actually visited by the evaluated policy at control step $t$.
% 其等权时间混合记为 $\bar\mu_\pi=T^{-1}\sum_{t=0}^{T-1}\mu_\pi^t$，
Its uniform mixture over time is $\bar\mu_\pi=T^{-1}\sum_{t=0}^{T-1}\mu_\pi^t$,
% 其中混合的样本保留时间索引。
with each sample retaining its time index.
% 有限时域性能差分恒等式给出
The finite-horizon performance-difference identity gives
\begin{equation}
 F(\pi)-F(\pi_0)
 =T\,\E_{\bar\mu_\pi}\!\left[A_t^0(x,\pi(x))\right].
 \label{eq:pd}
\end{equation}
% 式\eqref{eq:pd}沿待评价策略 $\pi$ 实际访问的情形计算动作后果，
Equation~\eqref{eq:pd} evaluates action effects at the situations actually visited by the evaluated policy $\pi$.
% 这一标准性能差分关系\cite{cheng2018value,ross2014aggrevate}
This standard performance-difference relation~\cite{cheng2018value,ross2014aggrevate}
% 为后续分析提供共同尺度，证明见附录~\ref{app:actionproofs}。
provides a common measure for the subsequent analysis; a proof is given in Appendix~\ref{app:actionproofs}.

% \subsubsection{动作误差何时能够约束失败概率}
\subsubsection{When Action Error Bounds Failure Probability}
\label{sec:actionerrorbound}
% 以专家策略 $\pi_{\mathrm E}$ 为参考，下面考察平均动作误差在什么条件下能够约束任务失败概率。
Taking the expert policy $\pi_{\mathrm E}$ as the reference, we ask when mean action error can bound task-failure probability.
% 以 $\nu$ 表示评价动作误差的参考分布，其样本是控制时刻与完整执行历史组成的 $(t,x_t)$。
% 相应的平均动作平方误差为 $R_\nu(\pi)=\E_{(t,x)\sim\nu}\norm{\pi(x)-\pi_{\mathrm E}(x)}^2$。
Let $\nu$ be the reference distribution used to evaluate action error, with samples $(t,x_t)$ specifying a control step and its complete execution history.
% 在专家示范场景中，$\nu$ 可取为示范采集过程诱导的总体分布，而不是已采集有限样本的经验分布。
For expert demonstrations, $\nu$ may be the population distribution induced by the demonstration-collection process, rather than the empirical distribution of the collected samples.
The corresponding mean squared action error is $R_\nu(\pi)=\E_{(t,x)\sim\nu}\norm{\pi(x)-\pi_{\mathrm E}(x)}^2$.
% $\nu$ 描述动作误差在哪里被衡量，而 $\bar\mu_\pi$ 描述策略在部署时实际遇到的情况。
Whereas $\nu$ specifies where action error is evaluated, $\bar\mu_\pi$ describes the situations actually encountered during deployment.
% 以下两个条件足以建立任务表现界。
The following two conditions suffice to establish a task-performance bound.

% \emph{覆盖：}任意一类执行情况在部署中出现的概率，不超过参考分布中对应概率的 $C$ 倍：
\emph{Coverage:} For any class of execution situations, its probability during deployment must be at most $C$ times its probability under the reference distribution:
\[
 \bar\mu_\pi(\mathcal B)\le C\,\nu(\mathcal B)
 \quad\text{for every measurable event }\mathcal B.
\]
% 其中，$\mathcal B$ 为时刻--历史空间中的一类执行情况，$1\le C<\infty$ 为统一常数。
% 该条件既排除参考分布完全未覆盖的部署情况，也限制参考分布对其出现概率的低估。
Here, $\mathcal B$ is a set of time--history pairs representing a class of execution situations, and $1\le C<\infty$ is a uniform constant.
This condition excludes deployment situations absent from the reference distribution and limits how much their probabilities can be underestimated.

% \emph{敏感性：}在固定动作尺度下，偏离专家动作所引起的失败概率变化必须受动作距离约束：
\emph{Sensitivity:} For a fixed action scale, the change in failure probability caused by deviating from the expert action must be bounded by the action distance:
\[
 |A_t^{\mathrm E}(x,a)|
 \le L\norm{a-\pi_{\mathrm E}(x)}.
\]
% 其中，$A_t^{\mathrm E}$ 为以专家作为参考策略时的动作影响，$0\le L<\infty$ 为在所考察时刻、历史及动作上统一成立的常数。
Here, $A_t^{\mathrm E}$ is the action effect defined using the expert as the reference policy,
and $0\le L<\infty$ is a uniform constant over the times, histories, and actions considered.
% 在这两个条件下，性能差分恒等式给出如下失败概率上界（证明见附录~\ref{app:actionproofs}）：
Under these conditions, the performance-difference identity yields the following bound (Appendix~\ref{app:actionproofs}):
\begin{equation}
 F(\pi)\le F(\pi_{\mathrm E})+TL\sqrt{C R_\nu(\pi)}.
 \label{eq:successbound}
\end{equation}
% 这一推论沿用由部署模仿误差约束任务代价的思路\cite{ross2011}，在此通过覆盖条件和连续动作敏感性连接参考动作误差与失败概率。
This corollary follows the use of deployment imitation error to bound task cost~\cite{ross2011}, here relating reference action error to failure probability through coverage and continuous-action sensitivity.
% 其中C控制参考误差向部署误差的转换，L控制动作偏差向任务后果的转换。
The factor $C$ controls the transfer from reference to deployment action error, while $L$ controls the conversion of action deviations into task effects.
% 在其他量不变且条件成立时，减小R、C或L均收紧上界，但不直接给出不同策略实际成功率的排序。
With the other quantities fixed and the conditions satisfied, reducing $R_\nu$, $C$, or $L$ tightens the bound, without establishing an ordering of actual policy success rates.
% 有限示范上的训练误差与总体误差R的联系还需单独的泛化论证。
Connecting finite-demonstration training error to the population error $R_\nu$ requires a separate generalization argument.
% L可能很大或缺少有限统一界，识别关键情形本身并不改变物理系统的敏感性。
The constant $L$ may be large or lack a finite uniform bound; identifying critical situations does not itself change physical sensitivity.

% \subsubsection{表示层面的覆盖与条件误差差异}
\subsubsection{Representation-Level Coverage and Conditional Error Discrepancy}
% 完整历史覆盖可因小幅动作偏差失效。
Coverage over complete histories can fail after a small action deviation (Appendix~\ref{app:historycoverage}).
% 将历史编码为任务表示可以合并无关差异，但也可能合并需要不同动作的情形。
Encoding histories into a task representation can merge irrelevant differences, but may also merge situations requiring different actions.

% 固定策略、专家和同一个保留时间索引的可测因果编码器Z；mu_Z和nu_Z是部署与参考分布诱导的表示分布。
For a fixed policy, expert, and common measurable causal encoder $Z=\phi(t,\mathcal H_t)$ retaining the time index, let $\mu_Z$ and $\nu_Z$ be the representation distributions induced by $\bar\mu_\pi$ and $\nu$.
% 表示覆盖要求每类表示在部署中的出现概率不超过参考概率的C_phi倍。
Suppose $\mu_Z(\mathcal B)\le C_\phi\nu_Z(\mathcal B)$ for every measurable event $\mathcal B$, with $1\le C_\phi<\infty$.
% 同一表示下，两种分布仍可能包含不同历史；以D_phi衡量其条件平均动作平方误差的差异。
Even at the same representation, the two distributions may contain different histories.
Let $e(x)=\norm{\pi(x)-\pi_{\mathrm E}(x)}$ and define their conditional-error discrepancy as $D_\phi=\E_{z\sim\mu_Z}|m_\mu(z)-m_\nu(z)|$,
where $m_\mu(z)=\E_{\bar\mu_\pi}[e(x)^2\mid Z=z]$ and $m_\nu(z)=\E_\nu[e(x)^2\mid Z=z]$.
% 当动作平方误差在两分布下可积、条件均值存在可测版本且B2的敏感性条件成立时，有如下风险迁移及任务界。
With square-integrable action error under both distributions, measurable versions of these conditional means, and the sensitivity condition of Section~\ref{sec:actionerrorbound},
\begin{align}
 \E_{\bar\mu_\pi}e(x)^2
 &\le C_\phi R_\nu(\pi)+D_\phi,
 \label{eq:representationrisktransfer-main}\\
 F(\pi)&\le F(\pi_{\mathrm E})
       +TL\sqrt{C_\phi R_\nu(\pi)+D_\phi}.
 \label{eq:representationcoveragebound}
\end{align}
% 第一式按表示条件化后迁移动作风险，第二式接入B2的敏感性分析；详细证明见附录。
The first inequality transfers action risk by conditioning on the representation; the second applies the sensitivity argument of Section~\ref{sec:actionerrorbound}.
Appendix~\ref{app:representationcoverage} gives the detailed proof.
% 该界同时取决于覆盖加权的参考风险与条件误差差异，覆盖常数较小本身不足以收紧上界。
The bound depends jointly on the coverage-weighted reference risk and the conditional error discrepancy; a smaller coverage factor alone need not tighten it.
% 保留RGB后附加确定性语义不自动改善覆盖。
Appending a deterministic semantic map to retained RGB does not by itself improve coverage.
% 对固定策略、专家与编码器，D_phi比较两分布的条件平均平方动作误差，而非任务信息损失；合并不同所需动作的历史也可能使其为零。
For the fixed policy, expert, and encoder, $D_\phi$ compares conditional mean squared action errors across distributions, rather than measuring lost task information.
It can vanish even when histories requiring different actions are merged, so a zero discrepancy does not establish information sufficiency.
% C直接以任务后果评价信息保留，与风险迁移分析互补；不界定D_phi，也不依赖上述覆盖条件。
Section~\ref{sec:representationanalysis} complements risk-transfer analysis by assessing information preservation directly through task consequences; it neither bounds $D_\phi$ nor requires the preceding coverage conditions.
